\section{Transformations of specifications}

Throughout, $(\Cfg, \mathcal{F}) = (E, \mathcal{E})^S$ is a countably infinite product,
$\Fin$ denotes the finite subsets of $S$, and $\lambda \in \Meas(E, \mathcal{E})$ is an a
priori measure.  Georgii §5.1 studies the invertible transformations of $\Cfg$ that are
induced by a bijection of the sites together with spin transformations acting separately at
each site, and their action on potentials, $\lambda$-modifications and specifications.

\begin{definition}[Transformations of $\Cfg$, Georgii (5.1)]
    \label{def:trans-transformation}
    \lean{MeasureTheory.GibbsMeasure.Transformation,
          MeasureTheory.GibbsMeasure.Transformation.toFun,
          MeasureTheory.GibbsMeasure.Transformation.measurable_toFun,
          MeasureTheory.GibbsMeasure.Transformation.toMeasurableEquiv,
          MeasureTheory.GibbsMeasure.Transformation.toMeasurableEquiv_eq_piCongr}
    \leanok

    Let $T$ be the set of all transformations $\tau : \Cfg \to \Cfg$ of the form
    \[\tau : \omega \mapsto (\tau_i\,\omega_{\tau_*^{-1} i})_{i \in S},\]
    where $\tau_*$ is a bijection of $S$ and each $\tau_i$ is a bimeasurable bijection of $E$.
    We write $\tau = (\tau_*; \tau_i, i \in S)$.  Every $\tau \in T$ is measurable, indeed a
    measurable isomorphism of $\Cfg$.

    In Lean, \texttt{Transformation S E} is the structure bundling the two data
    \texttt{sites : S $\simeq$ S} and \texttt{spin : S $\to$ E $\simeq^m$ E}; the action is
    \texttt{toFun}, and \texttt{toMeasurableEquiv} packages it as a measurable equivalence.
    It is identified with the composite of Mathlib's \texttt{piCongrLeft} (reindexing along
    $\tau_*$) and \texttt{piCongrRight} (applying the spins $\tau_i$).
\end{definition}

\begin{lemma}[$T$ is a group, Georgii §5.1]
    \label{lem:trans-group}
    \uses{def:trans-transformation}
    \lean{MeasureTheory.GibbsMeasure.Transformation.inv,
          MeasureTheory.GibbsMeasure.Transformation.comp,
          MeasureTheory.GibbsMeasure.Transformation.id,
          MeasureTheory.GibbsMeasure.Transformation.mul_def,
          MeasureTheory.GibbsMeasure.Transformation.one_def,
          MeasureTheory.GibbsMeasure.Transformation.inv_def,
          MeasureTheory.GibbsMeasure.Transformation.smul_def}
    \leanok

    $T$ is a group under composition, the inverse of $\tau = (\tau_*; \tau_i, i \in S)$ being
    $\tau^{-1} = (\tau_*^{-1}; \tau_{\tau_* i}^{-1}, i \in S)$, and $T$ acts on $\Cfg$ by
    $\tau \cdot \omega = \tau\omega$.
\end{lemma}
\begin{proof}
    \leanok

    Composition of the site bijections and of the spin equivalences is again of the displayed
    shape, and the stated formula for $\tau^{-1}$ is checked pointwise.  The Lean file supplies
    the \texttt{Group} instance on \texttt{Transformation S E} together with the
    \texttt{MulAction} on $S \to E$; the cited \texttt{mul\_def}, \texttt{one\_def},
    \texttt{inv\_def}, \texttt{smul\_def} identify the group operations with
    \texttt{comp}, \texttt{id}, \texttt{inv} and the action with \texttt{toFun}.
\end{proof}

\begin{lemma}[Transport of the finite-volume $\sigma$-algebras, Georgii, remark after (5.1)]
    \label{lem:trans-cylinder-transport}
    \uses{def:trans-transformation, def:cylinder-event}
    \lean{MeasureTheory.GibbsMeasure.Transformation.measurable_toFun_cylinderEvents,
          MeasureTheory.GibbsMeasure.Transformation.measurable_comp_cylinderEvents,
          MeasureTheory.GibbsMeasure.Transformation.measurable_toFun_cylinderEvents_compl,
          MeasureTheory.GibbsMeasure.Transformation.measurable_inv_toFun_cylinderEvents_compl}
    \leanok

    Let $\tau \in T$ and $\Lambda \subset S$.  Then $\tau$ is measurable from
    $\mathcal{F}_{\tau_*^{-1}\Lambda}$ to $\mathcal{F}_\Lambda$; consequently, if $f$ is
    $\mathcal{F}_\Lambda$-measurable then $f \circ \tau$ is
    $\mathcal{F}_{\tau_*^{-1}\Lambda}$-measurable.  The same holds for the complements
    $\Lambda^c$ and $(\tau_*^{-1}\Lambda)^c$, for $\tau$ and for $\tau^{-1}$.
\end{lemma}
\begin{proof}
    \leanok

    $\mathcal{F}_\Lambda$ is generated by the coordinate maps $\omega \mapsto \omega_j$,
    $j \in \Lambda$, and $(\tau\omega)_j = \tau_j\,\omega_{\tau_*^{-1}j}$ with
    $\tau_*^{-1} j \in \tau_*^{-1}\Lambda$; since $\tau_j$ is measurable, each such coordinate
    of $\tau$ is $\mathcal{F}_{\tau_*^{-1}\Lambda}$-measurable.  The complement statements
    follow because $\tau_*^{-1}(\Lambda^c) = (\tau_*^{-1}\Lambda)^c$.
\end{proof}

\begin{example}[The shift on $\mathbb{Z}^d$, Georgii (5.2)(1)]
    \label{ex:trans-shift}
    \uses{def:trans-transformation}
    \lean{MeasureTheory.GibbsMeasure.shift,
          MeasureTheory.GibbsMeasure.shift_toFun_apply,
          MeasureTheory.GibbsMeasure.shift_inv_toFun_apply}
    \leanok

    Let $S = \mathbb{Z}^d$.  For $j \in S$ the \emph{shift} (translation) by $j$ is
    \[\theta_j : \omega \mapsto (\omega_{i-j})_{i \in S},\]
    i.e.\ $\theta_j = (\,\cdot + j\,;\, \mathrm{id}_E, i \in S) \in T$.  The family
    $\Theta = (\theta_j)_{j \in S}$ is an abelian subgroup of $T$.

    In Lean the lattice $\mathbb{Z}^d$ is modelled as \texttt{Fin d $\to$ $\mathbb{Z}$} and
    $\theta_j$ is \texttt{shift E j}, with site part \texttt{Equiv.addRight j} and trivial
    spins; the two simp lemmas record $(\theta_j\omega)_i = \omega_{i-j}$ and
    $(\theta_j^{-1}\omega)_i = \omega_{i+j}$.  That $\Theta$ is an abelian subgroup is not
    recorded as a separate declaration.
\end{example}

\begin{example}[Spin flip and spin rotation, Georgii (5.2)(2),(3)]
    \label{ex:trans-spin-flip}
    \uses{def:trans-transformation}

    If $E$ is a symmetric Borel subset of $\R$, the \emph{spin flip}
    $\tau : \omega \mapsto (-\omega_i)_{i \in S}$ lies in $T$ (with $\tau_* = \mathrm{id}$).
    If $E$ is a rotation invariant Borel subset of $\R^N$ and $M \in SO(N)$, the \emph{spin
    rotation} $\omega \mapsto (M\omega_i)_{i \in S}$ lies in $T$.

    Both are instances of Definition~\ref{def:trans-transformation} with
    $\tau_* = \mathrm{id}$ and a constant family of spin equivalences.  Neither is formalized at
    Georgii's generality.  What the library has is the single Boolean instance
    \texttt{MeasureTheory.GibbsMeasure.Peierls.spinFlip}, the spin flip
    $\omega \mapsto (\neg\,\omega_i)_i$ on $\{0,1\}^{\mathbb{Z}^2}$ used for the Ising model in
    \texttt{Model/PhaseTransition.lean}; it is the case $E = \{\pm 1\}$ of (5.2)(2).  The
    general construction for a symmetric Borel $E \subset \R$, and the spin rotation of
    (5.2)(3), are not formalized.
\end{example}

\begin{definition}[$\lambda$-preserving transformations, Georgii §5.1]
    \label{def:trans-lambda-preserving}
    \uses{def:trans-transformation}

    Given $\lambda \in \Meas(E, \mathcal{E})$, a transformation
    $\tau = (\tau_*; \tau_i, i \in S) \in T$ is \emph{$\lambda$-preserving} if
    $\tau_i(\lambda) = \lambda \circ \tau_i^{-1} = \lambda$ for every $i \in S$.

    There is no named predicate for this in the library: it is carried as the hypothesis
    \texttt{h$\tau$ : $\forall$ i, MeasurePreserving ($\tau$.spin i) $\nu$ $\nu$} on the
    results that need it.
\end{definition}

\subsection*{The action on potentials, on modifications and on specifications}

\begin{definition}[Image of a potential, Georgii (5.3)]
    \label{def:trans-potential-map}
    \uses{def:trans-transformation}
    \lean{Potential.map, Potential.map_apply}
    \leanok

    For a family $\varphi = (\varphi_\Lambda)_{\Lambda \in \Fin}$ of functions on $\Cfg$ set
    \[\tau(\varphi) = (\varphi_{\tau_*^{-1}\Lambda} \circ \tau^{-1})_{\Lambda \in \Fin}.\]
    This applies in particular to potentials and to $\lambda$-modifications.  If $\Phi$ is a
    potential then so is $\tau(\Phi)$: the Lean file supplies the \texttt{IsPotential} instance
    for \texttt{Potential.map $\tau$ $\Phi$}, using
    Lemma~\ref{lem:trans-cylinder-transport}.
\end{definition}

\begin{lemma}[(5.3) is a left action]
    \label{lem:trans-potential-map-action}
    \uses{def:trans-potential-map, lem:trans-group}
    \lean{Potential.map_map, Potential.map_id, Potential.map_zero, Potential.map_add,
          Potential.map_smul}
    \leanok

    $\tau(\sigma(\Phi)) = (\tau \circ \sigma)(\Phi)$ and $\mathrm{id}(\Phi) = \Phi$; moreover
    $\Phi \mapsto \tau(\Phi)$ is $\R$-linear.
\end{lemma}
\begin{proof}
    \leanok

    Unfold (5.3) twice and use $(\tau\sigma)^{-1} = \sigma^{-1}\tau^{-1}$ together with
    $(\tau\sigma)_*^{-1}\Lambda = \sigma_*^{-1}\tau_*^{-1}\Lambda$.  Linearity is definitional,
    since (5.3) only reindexes and precomposes.
\end{proof}

\begin{definition}[Image of a specification, Georgii (5.4)]
    \label{def:trans-spec-map}
    \uses{def:trans-transformation, def:spec, def:proper-spec, def:markov-spec,
          lem:trans-cylinder-transport}
    \lean{Specification.map, Specification.map_apply, Specification.map_apply',
          Specification.map_apply_image}
    \leanok

    For a family $\gamma = (\gamma_\Lambda)_{\Lambda \in \Fin}$ of measure kernels on
    $\mathcal{F}$ the $\tau$-image is
    \[\tau(\gamma)_\Lambda(A \mid \omega) = \gamma_{\tau_*^{-1}\Lambda}(\tau^{-1}A \mid
      \tau^{-1}\omega) \qquad (\Lambda \in \Fin,\ A \in \mathcal{F},\ \omega \in \Cfg),\]
    equivalently $\tau(\gamma)_{\tau_*\Lambda}(\tau A \mid \tau\omega) =
    \gamma_\Lambda(A \mid \omega)$.  The $\tau$-image of a specification is a specification.

    In Lean \texttt{Specification.map} produces a bundled \texttt{Specification}, so the
    consistency, properness and Markov property of $\tau(\gamma)$ are part of the definition;
    \texttt{map\_apply'} and \texttt{map\_apply\_image} are the two displayed forms.
\end{definition}

\begin{lemma}[Georgii (5.5)]
    \label{lem:trans-spec-map-integral}
    \uses{def:trans-spec-map}
    \lean{Specification.lintegral_map_comp, Specification.integral_map_comp}
    \leanok

    Equation (5.4) is equivalent to
    \[(\tau(\gamma)_{\tau_*\Lambda} f) \circ \tau = \gamma_\Lambda(f \circ \tau)\]
    for all $\Lambda \in \Fin$ and all bounded or non-negative measurable $f$ on $\Cfg$.  The
    library proves both the Lebesgue form (for $f$ with values in $[0,\infty]$) and the
    Bochner form (for $f$ with values in a Banach space).
\end{lemma}
\begin{proof}
    \leanok

    By (5.4), $\tau(\gamma)_{\tau_*\Lambda}(\cdot \mid \tau\omega)$ is the image of
    $\gamma_\Lambda(\cdot\mid\omega)$ under $\tau$, so the identity is the change of variables
    formula along the measurable equivalence $\tau$.
\end{proof}

\begin{lemma}[(5.4) is a left action]
    \label{lem:trans-spec-map-action}
    \uses{def:trans-spec-map, lem:trans-group}
    \lean{Specification.map_id, Specification.map_comp}
    \leanok

    $\mathrm{id}(\gamma) = \gamma$ and $(\tau \circ \sigma)(\gamma) = \tau(\sigma(\gamma))$.
\end{lemma}
\begin{proof}
    \leanok

    Both sides are computed by \texttt{map\_apply}; the claim reduces to
    $(\tau\sigma)_*^{-1}\Lambda = \sigma_*^{-1}\tau_*^{-1}\Lambda$,
    $(\tau\sigma)^{-1} = \sigma^{-1}\tau^{-1}$ and functoriality of the pushforward of measures.
\end{proof}

\begin{definition}[Invariance and symmetries, Georgii (5.7)]
    \label{def:trans-invariant}
    \uses{def:trans-potential-map, def:trans-spec-map}
    \lean{Specification.IsInvariant, Specification.isInvariant_iff, Potential.map_eq_iff}
    \leanok

    Let $\tau \in T$ and $I \subset T$.
    \begin{enumerate}
        \item A family $\varphi = (\varphi_\Lambda)_{\Lambda \in \Fin}$ is \emph{$\tau$-invariant},
        and $\tau$ is a \emph{symmetry} of $\varphi$, if $\tau(\varphi) = \varphi$, i.e.\
        $\varphi_{\tau_*\Lambda} \circ \tau = \varphi_\Lambda$ for all $\Lambda \in \Fin$.
        \item A measure $\mu$ on $(\Cfg, \mathcal{F})$ is \emph{$\tau$-invariant} if
        $\tau(\mu) = \mu \circ \tau^{-1} = \mu$.  A specification $\gamma$ is
        \emph{$\tau$-invariant} if $\tau(\gamma) = \gamma$, i.e.\
        $\gamma_{\tau_*\Lambda}(\cdot \mid \tau\omega) = \tau(\gamma_\Lambda(\cdot\mid\omega))$
        for all $\Lambda \in \Fin$, $\omega \in \Cfg$.
        \item $\varphi$ (resp.\ $\mu$, $\gamma$) is \emph{$I$-invariant} if it is
        $\tau$-invariant for every $\tau \in I$.
    \end{enumerate}
    In Lean, $\tau$-invariance of a specification is \texttt{Specification.IsInvariant}, with
    \texttt{isInvariant\_iff} the displayed form of (b); \texttt{Potential.map\_eq\_iff} is the
    displayed form of (a) for potentials.  $\tau$-invariance of a measure is expressed by
    Mathlib's \texttt{MeasurePreserving $\tau$.toFun $\mu$ $\mu$}.
\end{definition}

\begin{lemma}[Symmetries form a monoid, Georgii (5.7)(d)]
    \label{lem:trans-symmetry-group}
    \uses{def:trans-invariant, lem:trans-spec-map-action}
    \lean{Specification.isInvariant_id, Specification.IsInvariant.comp}
    \leanok

    $\mathrm{id}$ is a symmetry of every specification $\gamma$, and if $\tau, \sigma$ are
    symmetries of $\gamma$ then so is $\tau \circ \sigma$.

    Georgii (5.7)(d) asserts more, namely that the symmetries of $\gamma$ form a \emph{subgroup}
    of $T$; closure under inverses is not proved in the tree, so what carries \leanok\ here is
    only the stated monoid version.
\end{lemma}
\begin{proof}
    \leanok

    Immediate from Lemma~\ref{lem:trans-spec-map-action}: $\mathrm{id}(\gamma) = \gamma$, and
    $(\tau\sigma)(\gamma) = \tau(\sigma(\gamma)) = \tau(\gamma) = \gamma$.
\end{proof}

\subsection*{Proposition (5.6): consistency of the three actions}

\begin{proposition}[Georgii (5.6)(a)]
    \label{prop:trans-isssd-invariant}
    \uses{def:trans-spec-map, def:trans-lambda-preserving, def:isssd, def:isssd-spec,
          def:juxtaposition, def:product-probability-measure}
    \lean{Specification.isssdFun_map_toFun, Specification.isssd_map_toFun,
          MeasureTheory.GibbsMeasure.Transformation.toFun_comp_juxt,
          MeasureTheory.GibbsMeasure.Transformation.measurePreserving_spin_piCongrLeft}
    \leanok

    Let $\lambda$ be an a priori measure and $\tau \in T$ be $\lambda$-preserving.  Then
    $\tau(\lambda_\cdot) = \lambda_\cdot$, i.e.\ the independent specification
    $\lambda_\cdot = \ISSSD(\lambda)$ is $\tau$-invariant.

    Georgii states this for $\lambda \in \Meas(E, \mathcal{E})$; the Lean statement assumes
    $\lambda$ is a probability measure (the standing hypothesis of \texttt{isssd}).  The
    underlying reindexing lemma \texttt{measurePreserving\_spin\_piCongrLeft} is proved for any
    $\sigma$-finite $\lambda$, hence is more general than the use made of it here.
\end{proposition}
\begin{proof}
    \leanok

    By definition $\lambda_\Lambda(\cdot\mid\omega)$ is the image of $\lambda^\Lambda$ under the
    juxtaposition $\zeta \mapsto \zeta_\Lambda \omega_{\Lambda^c}$.  The key identity
    \texttt{toFun\_comp\_juxt} says that $\tau$ intertwines the juxtapositions:
    $\tau \circ \juxtBy{\tau_*^{-1}\Lambda}(\tau^{-1}\omega)
     = \juxtBy{\Lambda}(\omega) \circ g$, where $g$ reindexes $\tau_*^{-1}\Lambda \to \Lambda$
    along $\tau_*$ and applies the spins $\tau_i$.  Since $\tau$ is $\lambda$-preserving, $g$
    carries $\lambda^{\tau_*^{-1}\Lambda}$ to $\lambda^{\Lambda}$, and the two pushforwards agree.
\end{proof}

\begin{proposition}[Georgii (5.6)(b)]
    \label{prop:trans-modification}
    \uses{prop:trans-isssd-invariant, def:modif, def:trans-potential-map, def:trans-spec-map}
    \lean{Specification.IsModifier.map_isssd, Specification.modification_isssd_map,
          Specification.coe_modification_isssd_map}
    \leanok

    Let $\tau \in T$ be $\lambda$-preserving.  For every $\lambda$-modification $\rho$, the
    family $\tau(\rho)$ of (5.3) is again a $\lambda$-modification, and
    \[\tau(\rho)\,\lambda_\cdot = \tau(\rho\,\lambda_\cdot).\]

    As in (5.6)(a), the Lean statement assumes that $\lambda$ is a probability measure.
\end{proposition}
\begin{proof}
    \leanok

    For bounded measurable $f$, (5.5) applied to $\gamma = \lambda_\cdot$ and (5.6)(a) give
    \[[\tau(\rho\lambda_\cdot)_{\tau_*\Lambda} f] \circ \tau
      = \lambda_\Lambda(\rho_\Lambda\, f \circ \tau)
      = \lambda_\Lambda\big((\tau(\rho)_{\tau_*\Lambda} f) \circ \tau\big)
      = [\lambda_{\tau_*\Lambda}(\tau(\rho)_{\tau_*\Lambda} f)] \circ \tau
      = [(\tau(\rho)\lambda_\cdot)_{\tau_*\Lambda} f] \circ \tau,\]
    and $\tau$ is invertible.  In Lean the computation is done at the level of kernels:
    \texttt{coe\_modification\_isssd\_map} rewrites the pushforward of a
    \texttt{withDensity} measure along the measurable equivalence $\tau$, and then applies
    Proposition~\ref{prop:trans-isssd-invariant}.  That $\tau(\rho)$ is a modification then
    follows because the resulting family of kernels is a specification.
\end{proof}

\begin{lemma}[Hamiltonians and partition functions, Georgii (5.6)(c)]
    \label{lem:trans-hamiltonian-map}
    \uses{def:trans-potential-map, lem:trans-spec-map-integral, prop:trans-isssd-invariant}
    \lean{Potential.hamiltonian_map, Potential.hamiltonian_map',
          Potential.boltzmannFactor_map, Specification.premodifierZ_map,
          Specification.premodifierNorm_map, Specification.lintegral_isssd_comp_inv}
    \leanok

    For every $\Lambda \in \Fin$,
    \[H^{\tau(\Phi)}_{\tau_*\Lambda} \circ \tau = H^\Phi_\Lambda, \qquad
      h^{\tau(\Phi)}_{\tau_*\Lambda} \circ \tau = h^\Phi_\Lambda,\]
    and, for $\lambda$-preserving $\tau$, $Z^{\tau(\rho)}_{\tau_*\Lambda} \circ \tau
    = Z^{\rho}_\Lambda$ for any family $\rho$ of densities, hence also
    $\tau(\rho)/Z^{\tau(\rho)} = \tau(\rho/Z^\rho)$.
\end{lemma}
\begin{proof}
    \leanok

    $H^{\tau(\Phi)}_{\tau_*\Lambda} \circ \tau
     = \sum_{A \cap \Lambda \neq \emptyset} \tau(\Phi)_{\tau_* A} \circ \tau
     = \sum_{A \cap \Lambda \neq \emptyset} \Phi_A = H^\Phi_\Lambda$, because
    $A \mapsto \tau_* A$ is a bijection of $\Fin$ preserving the summation filter and, by (5.3),
    $\tau(\Phi)_{\tau_* A} \circ \tau = \Phi_A$.  Exponentiating gives the statement for the
    Boltzmann factors, and applying (5.5) to $\gamma = \lambda_\cdot = \tau(\lambda_\cdot)$
    gives $Z^{\tau(\rho)}_{\tau_*\Lambda} \circ \tau
    = \lambda_\Lambda(h^{\tau(\Phi)}_{\tau_*\Lambda} \circ \tau) = Z^\Phi_\Lambda$.
\end{proof}

\begin{proposition}[Georgii (5.6)(c)]
    \label{prop:trans-admissible-map}
    \uses{lem:trans-hamiltonian-map, prop:trans-modification}
    \lean{Specification.IsPremodifierAdmissible.map,
          Potential.boltzmannFactor_map',
          Potential.premodifierNorm_boltzmannFactor_map}
    \leanok

    Let $\tau \in T$ be $\lambda$-preserving.  If $\Phi$ is a $\lambda$-admissible potential
    then $\tau(\Phi)$ is $\lambda$-admissible, and
    \[\rho^{\tau(\Phi)} = \tau(\rho^{\Phi}).\]
\end{proposition}
\begin{proof}
    \leanok

    Admissibility is $0 < Z_\Lambda < \infty$, which by
    Lemma~\ref{lem:trans-hamiltonian-map} is transported verbatim from $\Lambda$ to
    $\tau_*\Lambda$.  Since $\rho^\Phi_\Lambda = h^\Phi_\Lambda / Z^\Phi_\Lambda$ by (2.8), the
    same lemma gives $\rho^{\tau(\Phi)}_{\tau_*\Lambda} \circ \tau = \rho^\Phi_\Lambda
    = \tau(\rho^\Phi)_{\tau_*\Lambda} \circ \tau$, and $\tau$ is invertible.
\end{proof}

\begin{lemma}[The potential norms are transported]
    \label{lem:trans-norm-map}
    \uses{def:trans-potential-map}
    \lean{Potential.normAt_map, Potential.sum_normAt_map, Potential.hamiltonianBound_map,
          Potential.map_mem_absolutelySummable}
    \leanok

    $\|\tau(\Phi)\|_{\tau_* i} = \|\Phi\|_i$ for every $i \in S$, hence
    $\sum_{j \in \tau_*\Lambda}\|\tau(\Phi)\|_j = \sum_{i \in \Lambda}\|\Phi\|_i$; in particular
    $\tau(\Bspace) = \Bspace$, i.e.\ the action (5.3) preserves absolute summability, and
    likewise summability and finite range.

    This is not a numbered statement of Georgii; it records that the norms (2.12) and the bound
    (2.14) are transported by (5.3), which is what makes the closedness statement of (5.8)
    below available for an arbitrary $\tau \in T$.
\end{lemma}
\begin{proof}
    \leanok

    $\|\Phi\|_i = \sum_{A \ni i}\sup_\eta |\Phi_A(\eta)|$; reindex the sum along the bijection
    $A \mapsto \tau_* A$ of $\Fin$, which matches $\{A : i \in A\}$ with
    $\{B : \tau_* i \in B\}$, and use that $\eta \mapsto \tau\eta$ is a bijection of $\Cfg$, so
    the suprema agree.
\end{proof}

\subsection*{Corollary (5.9) and shift-invariance}

\begin{corollary}[Georgii (5.9)(b), forward direction]
    \label{cor:trans-gibbs-spec}
    \uses{prop:trans-admissible-map, def:trans-invariant, def:gibbs-meas}
    \lean{Potential.map_gibbsSpecification, Potential.isInvariant_gibbsSpecification,
          Potential.gibbsSpecification_congr}
    \leanok

    Let $\tau \in T$ be $\lambda$-preserving and $\Phi$ a $\lambda$-admissible potential.  Then
    $\tau(\gamma^\Phi) = \gamma^{\tau(\Phi)}$; in particular, if $\Phi$ is $\tau$-invariant then
    so is $\gamma^\Phi$.

    The Lean statements are at \emph{stronger} hypotheses than Georgii's: they require $S$
    countable and $\Phi \in \Bspace$ absolutely summable (the setting in which
    \texttt{gibbsSpecificationOfAbsolutelySummable} is defined), whereas Georgii only assumes
    $\Phi$ to be $\lambda$-admissible.  Proposition~\ref{prop:trans-admissible-map}, which is
    the substance of the argument, is available at Georgii's generality.
\end{corollary}
\begin{proof}
    \leanok

    $\gamma^\Phi = \rho^\Phi \lambda_\cdot$ by definition.  By (5.6)(b),
    $\tau(\rho^\Phi \lambda_\cdot) = \tau(\rho^\Phi)\lambda_\cdot$, and by (5.6)(c)
    $\tau(\rho^\Phi) = \rho^{\tau(\Phi)}$, whence
    $\tau(\gamma^\Phi) = \rho^{\tau(\Phi)}\lambda_\cdot = \gamma^{\tau(\Phi)}$.  If
    $\tau(\Phi) = \Phi$ this reads $\tau(\gamma^\Phi) = \gamma^\Phi$.
\end{proof}

\begin{remark}[The converses in Georgii (5.9)]
    \label{rmk:trans-59-converse}
    \uses{cor:trans-gibbs-spec, def:trans-invariant}

    Georgii (5.9) also contains two converses, neither of which is formalized.
    \begin{enumerate}
        \item (5.9)(a): if $E$ is countable, $\mathcal{E} = \mathfrak{P}(E)$, $\lambda$ is
        equivalent to counting measure and $\gamma$ is a $\tau$-invariant specification, then the
        unique $\lambda$-modification $\rho$ with $\gamma = \rho\lambda_\cdot$ is
        $\tau$-invariant.  The proof uses the explicit formula
        $\rho_\Lambda(\omega) \propto \gamma_\Lambda(\sigma_\Lambda = \omega_\Lambda \mid \omega)$
        of Remark (1.28), which the library records only informally.
        \item (5.9)(b), converse: if $\Phi$ is normalized by some $\alpha \in \Prob(E, \mathcal{E})$
        preserved by $\tau$, and either $\alpha$ is a Dirac measure or $\Phi$ converges uniformly,
        then $\tau$-invariance of $\rho^\Phi$ implies $\tau$-invariance of $\Phi$.  This needs
        the equivalence theorems (2.34) and (2.35)(a) for potentials, which are not in the tree.
    \end{enumerate}
    Consequently the concluding statement of §5.1 --- for $S = \mathbb{Z}^d$, a normalized
    potential $\Phi$ is shift-invariant \emph{iff} $\rho^\Phi$ is --- is available in the library
    in one direction only.
\end{remark}

\begin{example}[Shift-invariance, Georgii (5.8)]
    \label{ex:trans-shift-invariant}
    \uses{ex:trans-shift, def:trans-invariant, def:trans-potential-map}
    \lean{Potential.IsShiftInvariant, Potential.isShiftInvariant_iff}
    \leanok

    Let $S = \mathbb{Z}^d$.  A specification $\gamma$ is \emph{shift-invariant} (homogeneous) if
    it is $\Theta$-invariant, i.e.\
    $\gamma_{\Lambda+j}(\theta_j A \mid \theta_j\omega) = \gamma_\Lambda(A\mid\omega)$ for all
    $\Lambda \in \Fin$, $j \in S$, $A \in \mathcal{F}$, $\omega \in \Cfg$.  A potential $\Phi$ is
    shift-invariant if $\Phi_{A+j} \circ \theta_j = \Phi_A$ for all $A \in \Fin$, $j \in S$.  We
    write $\Bspace_\Theta$ for the set of shift-invariant elements of $\Bspace$.

    The library defines \texttt{Potential.IsShiftInvariant} as
    $\forall j,\ \theta_j(\Phi) = \Phi$ and \texttt{isShiftInvariant\_iff} unfolds it to
    Georgii's displayed form.  Shift-invariance of a \emph{specification} is not given a
    separate name: it is $\forall j$, \texttt{Specification.IsInvariant (shift E j) $\gamma$}.
\end{example}

\begin{theorem}[$\Bspace_\Theta$ is closed, Georgii (5.8)]
    \label{thm:trans-b-theta-closed}
    \uses{ex:trans-shift-invariant, lem:trans-norm-map}
    \lean{Potential.isClosed_setOf_map_eq, Potential.isClosed_setOf_isShiftInvariant}
    \leanok

    $\Bspace_\Theta$ is a $\|\cdot\|_0$-closed subspace of $\Bspace$.

    The library proves \emph{more}: for \emph{any} $\tau \in T$ (and $S$ countable), the fixed
    point set $\{\Phi \in \Bspace : \tau(\Phi) = \Phi\}$ is closed in $\Bspace$; Georgii's
    statement is the intersection of these sets over the shift group $\Theta$.  The extra
    hypothesis $S$ countable is automatic for $S = \mathbb{Z}^d$.
\end{theorem}
\begin{proof}
    \leanok

    The fixed point set is characterised by
    $\Phi_{\tau_* A}(\tau\eta) = \Phi_A(\eta)$ for all $A \in \Fin$, $\eta \in \Cfg$
    (Definition~\ref{def:trans-invariant}), i.e.\ as a set of coincidences of two evaluation
    maps.  On $\Bspace$ the norm $\|\cdot\|_0$ dominates each such evaluation, so these maps are
    continuous and the set where they agree is (sequentially) closed.  Intersecting over
    $j \in \mathbb{Z}^d$ gives $\Bspace_\Theta$.
\end{proof}

\begin{remark}[What (5.8) states beyond closedness]
    \label{rmk:trans-b-theta-banach}
    \uses{thm:trans-b-theta-closed}

    Georgii (5.8) additionally asserts that $\Bspace_\Theta$ is a Banach space for the norm
    $\|\Phi\|_0 = \sum_{A \ni 0} \|\Phi_A\|$, and that the shift-invariant potentials of bounded
    range are dense in $\Bspace_\Theta$.  Neither is in the library: what is available is the
    closedness of Theorem~\ref{thm:trans-b-theta-closed} inside the (complete) space $\Bspace$,
    and Lemma~\ref{lem:trans-norm-map}, which shows $\|\cdot\|_0$ is shift-invariant.  Not yet
    formalized.
\end{remark}

\begin{remark}[Outlook: Gibbs measures with symmetries]
    \label{rmk:trans-outlook}
    \uses{def:trans-invariant, def:gibbs-meas, def:trans-spec-map}

    The action (5.4) transports Gibbs measures: $\mu \in \Gibbs{\gamma}$ implies
    $\tau(\mu) \in \Gibbs{\tau(\gamma)}$, so $\Gibbs{\gamma}$ is invariant under every symmetry
    of $\gamma$, and if $\Gibbs{\gamma} = \{\mu\}$ then $\mu$ inherits all symmetries of
    $\gamma$.  These are Georgii (5.10) and (5.11) and belong to §5.2; the library has them as
    \texttt{Specification.map\_mem\_GP}, \texttt{Specification.IsInvariant.map\_mem\_GP},
    \texttt{Specification.IsInvariant.map\_isGibbsMeasure} and
    \texttt{Specification.IsInvariant.measurePreserving\_of\_GP\_eq\_singleton}, and the
    $L$-closedness of the $I$-invariant random fields and of $\mathcal{G}_I(\gamma)$ as
    \texttt{MeasureTheory.GibbsMeasure.isClosed\_setOf\_forall\_measurePreserving} and
    \texttt{MeasureTheory.GibbsMeasure.isClosed\_setOf\_mem\_GP\_and\_measurePreserving}.  They
    are stated and proved in the next section.
\end{remark}
